% C090 | 椭圆与双曲线共焦点的焦角关系（精简源稿；依赖项目LaTeX模板）
\textbf{C090\quad 共焦椭圆、双曲线焦角与离心率关系}

设非退化椭圆与双曲线有共同焦点 $F_1,F_2$，且交于点 $P$；
$e_1\in(0,1)$、$e_2>1$ 分别为椭圆和双曲线的离心率，
$\theta=\angle F_1PF_2\in(0,\pi)$，则
\[
\boxed{\frac{1-\cos\theta}{e_1^2}
+\frac{1+\cos\theta}{e_2^2}=2}
\quad\Longleftrightarrow\quad
\boxed{\frac{\sin^2(\theta/2)}{e_1^2}
+\frac{\cos^2(\theta/2)}{e_2^2}=1}.
\]
推论（柯西不等式）：
\[
\boxed{\frac1{e_1}+\frac1{e_2}\le\frac2{\sin\theta}}.
\]
等号成立当且仅当 $e_1=\tan(\theta/2)$、
$e_2=\cot(\theta/2)$；在椭圆、双曲线的离心率限制下，
等号只可能出现在 $0<\theta<\pi/2$。
\textbf{最短主线：}同一焦点三角形的「和」来自椭圆、「差」来自双曲线，
余弦定理直接连接焦角与离心率。
